\section{本章小结}
\label{sec:dsel-summary}

本章把概率分布当作未知量，在允许的分布中按目标与约束选出一个。章首三种愿望的分歧有了明确的回答：降低期望成本是要最小化的线性目标，单独使用会把质量推到顶点；指定平均量是等式约束，单独使用留下无穷多个候选；接近参考分布是以KL计量的代价，单独使用什么也不改变。把它们组合成一个问题，再写清候选集合中的非负性、归一化、支持与可积性，问题才有确定的含义；第一节的可数例子说明，省略可积性会让目标变成依赖配对方式的“$-\infty+\infty$”。

三类问题的解都是“参考测度乘以指数权重再归一化”，但各自的来源与条件不同，表~\ref{tab:dsel-summary}把它们并列。

\begin{BookInlineTable}
 \centering
 \caption{本章三类分布选择问题。$A(\bm{\theta})$是以参考测度为基准的对数配分函数，$Z=\mathbb{E}_{P}[\mathrm{e}^{-V/\tau}]$。矩约束最小KL的条件指有限状态、参考概率处处为正的情形；最大熵还要求候选分布的熵有定义。}
 \label{tab:dsel-summary}
 \small
 \begin{tabular}{llll}
 \toprule
 问题 & 解的形式 & 最优值 & 主要条件 \\
 \midrule
 KL正则化成本 & $\mathrm{d}Q^{*}/\mathrm{d}P\propto\mathrm{e}^{-V/\tau}$ & $-\tau\ln Z$ & $Z<\infty$，$V\in L^{1}(Q^{*})$ \\
 矩约束最小KL & $\mathrm{d}Q^{*}/\mathrm{d}P\propto\mathrm{e}^{\langle\bm{\theta},\bm{T}\rangle}$ & $\langle\bm{\theta},\bm{m}\rangle-A(\bm{\theta})$ & $\bm{m}\in\operatorname{ri}\mathcal{M}$ \\
 最大熵 & $\mathrm{d}Q^{*}/\mathrm{d}\nu\propto\mathrm{e}^{\langle\bm{\theta},\bm{T}\rangle}$ & $A(\bm{\theta})-\langle\bm{\theta},\bm{m}\rangle$ & 配分积分有限 \\
 \bottomrule
 \end{tabular}
\end{BookInlineTable}

Gibbs分布由KL正则化的成本问题给出，它的最优性由恒等式~\eqref{eq:dsel-gibbs-identity}证明：任何候选比最优值多出的部分恰是$\tau D_{\mathrm{KL}}(Q\Vert Q^{*})$。指数倾斜由矩约束下的最小KL给出，最优性由勾股恒等式证明，求解归结为$k$个未知数的矩匹配方程；方程的Jacobian是统计量的协方差，冗余统计量使参数不唯一而分布仍唯一，边界上的矩使最优分布收缩支持、参数发散。最大熵在有限支持上就是以均匀分布为参考的最小KL；在连续情形中熵依赖于参考测度，只有参考密度的对数是约束统计量的仿射函数时，最大熵才与最小KL重合。

这三种解的指数形式都是由最优性推出来的，不是预先选定的分布族。公式只给出候选，第五节的检查把“候选合法”“最小值存在”“候选达到最小值”分开处理：对偶下界与可行上界之差给出可计算的最优性证书；无穷状态上最小值可能达不到，写得出的Gibbs分布也可能因可积性失败而不是候选。

本章的出发点是参考分布与约束，目标分布并不存在，要选的就是它。下一章换一个起点：目标分布已经给定，但难以计算，我们限制候选集合，在容易表示和使用的分布中寻找它的近似。那时KL的两个方向、候选族的表达能力与逼近误差，成为需要分别研究的问题。
